% !TeX program = xelatex
% !TeX encoding = UTF-8
% 编译方法：xelatex 论文模版.tex（建议连续编译两次以更新交叉引用）
\documentclass[UTF8,a4paper,zihao=-4,fontset=none]{ctexart}

% 字体采用三级自动选择机制，使模板能够在不同电脑上直接编译。
% 先检查 Windows 标准字体目录；再检查 macOS 版 Microsoft Word 的字体；均不可用时回退到 TeX Live 自带的 Fandol 字体。
% 无论选择哪一组字体，主字体的 BoldFont 都显式绑定到对应黑体，确保 \textbf 下的中文字体正确。
\newcommand{\UseWindowsMicrosoftFonts}{
  \setCJKmainfont{simsun.ttc}[Path={C:/Windows/Fonts/},BoldFont=simhei.ttf]
  \setCJKsansfont{simhei.ttf}[Path={C:/Windows/Fonts/},BoldFont=simhei.ttf]
  \setCJKmonofont{simsun.ttc}[Path={C:/Windows/Fonts/}]
  \setCJKfamilyfont{paperhei}{simhei.ttf}[Path={C:/Windows/Fonts/},BoldFont=simhei.ttf]
  \PackageInfo{paper-template}{Using Windows SimSun and SimHei}
}

\newcommand{\UseOfficeMicrosoftFonts}{
  \setCJKmainfont{Simsun.ttc}[
    Path={/Applications/Microsoft Word.app/Contents/Resources/DFonts/},
    BoldFont=SimHei.ttf
  ]
  \setCJKsansfont{SimHei.ttf}[
    Path={/Applications/Microsoft Word.app/Contents/Resources/DFonts/},
    BoldFont=SimHei.ttf
  ]
  \setCJKmonofont{Simsun.ttc}[
    Path={/Applications/Microsoft Word.app/Contents/Resources/DFonts/}
  ]
  \setCJKfamilyfont{paperhei}{SimHei.ttf}[
    Path={/Applications/Microsoft Word.app/Contents/Resources/DFonts/},
    BoldFont=SimHei.ttf
  ]
  \PackageInfo{paper-template}{Using Microsoft Office SimSun and SimHei}
}

\newcommand{\UsePortableFandolFonts}{
  \setCJKmainfont{FandolSong-Regular.otf}[cmap=UniGB-UTF16-H,BoldFont=FandolHei-Regular.otf]
  \setCJKsansfont{FandolHei-Regular.otf}[cmap=UniGB-UTF16-H,BoldFont=FandolHei-Regular.otf]
  \setCJKmonofont{FandolFang-Regular.otf}[cmap=UniGB-UTF16-H]
  \setCJKfamilyfont{paperhei}{FandolHei-Regular.otf}[cmap=UniGB-UTF16-H,BoldFont=FandolHei-Regular.otf]
  \PackageWarningNoLine{paper-template}{SimSun and SimHei not found; using portable Fandol fonts}
}

\newcommand{\UseOfficeOrPortableFonts}{
  \IfFileExists{/Applications/Microsoft Word.app/Contents/Resources/DFonts/Simsun.ttc}{
    \IfFileExists{/Applications/Microsoft Word.app/Contents/Resources/DFonts/SimHei.ttf}{
      \UseOfficeMicrosoftFonts
    }{
      \UsePortableFandolFonts
    }
  }{
    \UsePortableFandolFonts
  }
}

% 将下面的 0 改为 1，可强制使用完全随 TeX Live 提供的 Fandol 字体进行兼容性测试。
\providecommand{\ForcePortableFonts}{0}
\ifnum\ForcePortableFonts=1
  \UsePortableFandolFonts
\else
  \IfFileExists{C:/Windows/Fonts/simsun.ttc}{
    \IfFileExists{C:/Windows/Fonts/simhei.ttf}{
      \UseWindowsMicrosoftFonts
    }{
      \UseOfficeOrPortableFonts
    }
  }{
    \UseOfficeOrPortableFonts
  }
\fi
\newcommand{\heifont}{\CJKfamily{paperhei}}

% 页面、行距与段落。
\usepackage[a4paper,top=2.5cm,bottom=2.5cm,left=2.5cm,right=2.5cm,footskip=1.2cm]{geometry}
\usepackage{setspace}
\singlespacing
\setlength{\parindent}{2em}
\setlength{\parskip}{0pt}

% 数学公式、表格与插图。
\usepackage{amsmath,amssymb,mathtools,bm}
\usepackage{graphicx}
\usepackage{booktabs,tabularx,array,multirow}
\usepackage[table]{xcolor}
\usepackage{float}
\usepackage{caption}
\usepackage{subcaption}
\usepackage{enumitem}

% 标题、页码与交叉引用。
\usepackage{titlesec}
\usepackage{fancyhdr}
\usepackage[hidelinks]{hyperref}
\usepackage{cleveref}

% 一级标题使用中文序号并居中，下级标题使用阿拉伯数字层级编号并左对齐。
\renewcommand{\thesection}{\arabic{section}}
\renewcommand{\thesubsection}{\arabic{section}.\arabic{subsection}}
\renewcommand{\thesubsubsection}{\arabic{section}.\arabic{subsection}.\arabic{subsubsection}}
\setcounter{secnumdepth}{3}
\titleformat{\section}[block]
  {\centering\heifont\bfseries\zihao{4}}
  {\chinese{section}、}{0pt}{}
\titleformat{\subsection}[hang]
  {\raggedright\heifont\bfseries\zihao{-4}}
  {\thesubsection}{0.8em}{}
\titleformat{\subsubsection}[hang]
  {\raggedright\heifont\bfseries\zihao{-4}}
  {\thesubsubsection}{0.8em}{}
\titlespacing*{\section}{0pt}{2.8ex plus 0.5ex minus 0.2ex}{1.7ex}
\titlespacing*{\subsection}{0pt}{1.6ex plus 0.3ex minus 0.1ex}{0.8ex}
\titlespacing*{\subsubsection}{0pt}{1.4ex plus 0.2ex minus 0.1ex}{0.7ex}

% 页码从摘要页开始，位于页脚中央。
\pagestyle{fancy}
\fancyhf{}
\fancyfoot[C]{\zihao{5}\thepage}
\renewcommand{\headrulewidth}{0pt}
\renewcommand{\footrulewidth}{0pt}

% 图表及列表的常用设置。
\captionsetup{font=small,labelfont=bf,labelsep=quad}
\setlist[itemize]{itemsep=0.35\baselineskip,topsep=0.8ex,parsep=0pt,partopsep=0pt}
\newcolumntype{C}[1]{>{\centering\arraybackslash}p{#1}}
\newcolumntype{L}[1]{>{\raggedright\arraybackslash}p{#1}}
\newcolumntype{Y}{>{\raggedright\arraybackslash}X}
\newcolumntype{Z}{>{\centering\arraybackslash}X}

% 请优先修改下面三个字段。
\newcommand{\papertitle}{论文题目}
\newcommand{\paperabstract}{请在此填写摘要。摘要应简要说明研究问题、所用方法、主要结果与结论，避免写成正文各节的简单罗列。}
\newcommand{\paperkeywords}{关键词一；关键词二；关键词三}

\hypersetup{
  pdftitle={\papertitle},
  pdfauthor={}
}

\begin{document}
\pagenumbering{arabic}
\setcounter{page}{1}
\thispagestyle{fancy}

% 第一页：标题、摘要与关键词。
\begin{center}
  \vspace*{0.3cm}
  {\heifont\bfseries\zihao{3}\papertitle\par}
  \vspace{1.5\baselineskip}
  {\heifont\bfseries\zihao{4}摘要\par}
\end{center}

\paperabstract\par
\vspace{0.5\baselineskip}
\noindent{\heifont\bfseries 关键词：}\paperkeywords

\newpage

\section{问题重述}

\subsection{问题背景}

请在此介绍问题产生的现实背景、研究对象与研究意义。对题目材料中的事实应准确转述，不宜在本小节提前展开模型推导。

\subsection{问题重述}

请在此用自己的语言准确概括需要解决的各项任务，明确已知条件、待求目标及各问之间的关系。

\section{模型假设}

为使模型可解且与实际问题相符，可在此列出必要假设，并逐条说明假设依据。
\begin{itemize}[label=\textbullet,leftmargin=2em]
  \item \textbf{假设一：}请在此填写假设内容及其依据。
  \item \textbf{假设二：}请在此填写假设内容及其依据。
  \item \textbf{假设三：}请在此填写假设内容及其依据。
\end{itemize}

\section{符号说明}

\begin{table}[H]
  \centering
  \begingroup
  \renewcommand{\arraystretch}{1.45}
  \begin{tabularx}{\textwidth}{C{0.18\textwidth}ZC{0.18\textwidth}}
    \toprule[1.2pt]
    {\heifont\bfseries 符号} & {\heifont\bfseries 定义} & {\heifont\bfseries 单位} \\
    \midrule[0.5pt]
    \(\Omega\) & 以原点为圆心、半径为 \(1800\,\mathrm{m}\) 的目标圆域 & -- \\
    \(g\) & 干扰源的位置，固定但未知 & \(\mathrm{m}\) \\
    \(s_i\) & 第 \(i\) 个检测点的位置 & \(\mathrm{m}\) \\
    \(q\) & 待选择的检测点或行动位置 & \(\mathrm{m}\) \\
    \(R\) & 干扰源的固定有效接收半径，取值为 \(1000\)--\(1500\,\mathrm{m}\) & \(\mathrm{m}\) \\
    \(\theta_i\) & 在检测点 \(s_i\) 获得的正常示向度读数 & \(\mathrm{rad}\) \\
    \(\alpha\) & 测角误差上界，\(\alpha=\pi/180\) & \(\mathrm{rad}\) \\
    \(r_{\mathrm{strong}}\) & 强信号距离阈值，取值为 \(5\,\mathrm{m}\) & \(\mathrm{m}\) \\
    \(r_{\mathrm{opt}}\) & 光学精确定位距离，取值为 \(20\,\mathrm{m}\) & \(\mathrm{m}\) \\
    \(\mathcal W_i\) & 第 \(i\) 次正常测向对应的正向角域，半角为 \(\alpha\) & -- \\
    \(P\) & 纯测角交会区域，即各角域 \(\mathcal W_i\) 的交集 & -- \\
    \(K\) & 综合观测信息与物理条件得到的源位置可行区域 & -- \\
    \(\operatorname{diam}(K)\) & 区域 \(K\) 的直径 & \(\mathrm{m}\) \\
    \(c(K)\) & 区域 \(K\) 的最小包围圆圆心 & \(\mathrm{m}\) \\
    \(\rho(K)\) & 区域 \(K\) 的最小包围圆半径 & \(\mathrm{m}\) \\
    \(d_{\max}(q;K)\) & 位置 \(q\) 到区域 \(K\) 的最远距离 & \(\mathrm{m}\) \\
    \bottomrule[1.2pt]
  \end{tabularx}
  \endgroup
\end{table}

上述位置变量均为二维坐标向量；公式推导中的角度统一采用弧度，观测数据与结果展示可采用度并注明单位。多源情形下增加源编号下标，其余局部符号在各问建模时另行定义。

\section{问题分析}

\subsection{问题一分析}

请在此说明问题一的核心矛盾、数据特征、可采用的方法以及选择该方法的理由。

\subsection{问题二分析}

请在此说明问题二与问题一之间的联系，以及需要新增或调整的模型要素。

\section{问题一的模型建立与求解}

\subsection{模型建立}

请在此给出变量关系、目标函数、约束条件和参数含义。数学公式可使用 equation、align 等环境，并通过 label 与 cref 命令交叉引用。

\subsection{模型求解}

请在此说明求解方法、计算过程与主要结果。

\section{问题二的模型建立与求解}

\subsection{模型建立}

请在此说明问题二所需的新变量、模型关系和约束条件。

\subsection{模型求解}

请在此说明问题二的求解过程与主要结果。其余问题可按相同层级继续添加新的 section。

\section{模型检验}

请在此通过误差分析、敏感性分析、稳健性检验或与实际数据对照等方式检验模型。

\section{模型评价与推广}

\subsection{模型的优点}

请在此概括模型在合理性、解释性、精度或计算效率方面的优点。

\subsection{模型的局限性}

请在此客观说明假设、数据和方法带来的限制。

\subsection{模型的推广}

请在此说明模型能够扩展到哪些情形，以及推广时需要修改的条件。

\begin{thebibliography}{99}
  \bibitem{reference1} 作者. 文献题目[J]. 期刊名称, 年份, 卷号(期号): 起止页码.
  \bibitem{reference2} 作者. 书名[M]. 出版地: 出版社, 年份.
\end{thebibliography}

\end{document}
